% ============================================================================
% Fichier  : partie4/P04_C16_api_mobile.tex
% Titre    : Interfaces de Programmation et Applications Mobiles
% Auteur   : MINKA MI NGUIDJOÏ Thierry Emmanuel
% Version  : 1.0 -- novembre 2026
% ============================================================================
% Sources : OWASP API Security Top 10:2023 ; OWASP Mobile Application
% Security Verification Standard (MASVS). Opération MEFOU (juillet 2026) :
% cas fictif.
% ============================================================================

\chapter{Interfaces de Programmation et Applications Mobiles}
\label{chap:api}
\chaptermark{Interfaces de programmation et mobiles}

\epigraph{%
    « L'application mobile n'est qu'une façade~;\\
    le vrai test porte sur l'interface qu'elle appelle. »%
}{Principe du cours}

\niveauChapitre{Intermédiaire à expert}{Chapitre~\ref{chap:injections}.}

\begin{boxobjectifs}
\begin{itemize}
    \item Décrire une interface de programmation (API) du point de vue du
          testeur.
    \item Reconnaître les risques du Top~10 des API de l'OWASP, à commencer
          par le défaut d'accès au niveau de l'objet.
    \item Comprendre pourquoi l'application mobile ne protège pas le serveur.
    \item Tester une application mobile selon le standard MASVS.
    \item Rédiger les constats d'une mission mobile.
\end{itemize}
\end{boxobjectifs}

\bigskip

% ============================================================================
\section*{Question directrice : l'application mobile de LiZbiZ est-elle plus sûre que son portail web~?}
% ============================================================================

En juillet, avant le test principal, LiZbiZ confie à NYONG SÉCURITÉ un test
de son application mobile de suivi, l'Opération MEFOU. L'équipe découvre
vite que l'application n'est qu'une fenêtre~: toute la logique est dans
l'interface de programmation qu'elle appelle, et c'est là que se trouvent
les faiblesses. Six constats en ressortent.

\medskip
\textbf{Plan du chapitre.}
\begin{enumerate}
    \item L'interface de programmation vue par le testeur
          (section~\ref{sec:c16-api}).
    \item Le Top~10 des API (section~\ref{sec:c16-top10}).
    \item La façade mobile (section~\ref{sec:c16-facade}).
    \item Tester le mobile avec le MASVS (section~\ref{sec:c16-masvs}).
    \item L'Opération MEFOU (section~\ref{sec:c16-mefou}).
\end{enumerate}

\separateur

% ============================================================================
\section{L'Interface de Programmation vue par le Testeur}
\label{sec:c16-api}
% ============================================================================

\begin{boxdef}
\termecle{Interface de programmation} (API)~: ensemble de points d'entrée
par lesquels des programmes échangent des données, le plus souvent en HTTP,
avec des messages structurés. Le testeur l'interroge directement, sans
passer par l'application qui s'en sert.
\end{boxdef}

Une interface de programmation reçoit souvent moins d'attention que le site
web~: pas d'écran, pas de visiteur humain, l'illusion qu'elle est « interne ».
C'est une erreur~: elle est aussi exposée, et porte la logique métier.

% ============================================================================
\section{Le Top 10 des API}
\label{sec:c16-top10}
% ============================================================================

\begin{tableaucours}{L{4.6cm}L{7.6cm}}{%
    \tableauHeader{Risque (OWASP API 2023)} &
    \tableauHeader{Question du testeur}
}
Accès au niveau de l'objet & Puis-je lire ou modifier l'objet d'un autre en
changeant son identifiant~? \\
Authentification défaillante & Les jetons sont-ils solides, expirent-ils~? \\
Accès au niveau d'une propriété & Puis-je lire ou écrire un champ qui ne
m'est pas destiné~? \\
Consommation de ressources & Puis-je épuiser le service par des appels
massifs~? \\
Accès au niveau de la fonction & Puis-je appeler une fonction réservée à
un administrateur~? \\
\end{tableaucours}
\vspace{-0.8em}
\begin{center}{\small\itshape Tableau~16.1. Principaux risques des API~\cite{owaspapitop10}}\end{center}

Le premier risque, le défaut d'accès au niveau de l'objet, est l'équivalent
de la référence directe du Chapitre~\ref{chap:injections}, appliqué aux
interfaces de programmation~: c'est le plus fréquent et le plus grave.

% ============================================================================
\section{La Façade Mobile}
\label{sec:c16-facade}
% ============================================================================

\begin{boxalert}
\textbf{Violation typique.} Croire qu'un contrôle présent dans
l'application mobile protège le service.

\smallskip
\textbf{Règle~:} tout ce que fait l'application mobile, un programme peut le
refaire en appelant directement l'interface, sans ses contrôles. Une clé
cachée dans l'application n'est pas secrète~: elle se lit dans le paquet
installé.
\end{boxalert}

% ============================================================================
\section{Tester le Mobile avec le MASVS}
\label{sec:c16-masvs}
% ============================================================================

Le \termecle{standard de vérification de la sécurité des applications
mobiles} (MASVS) de l'OWASP fixe les exigences à vérifier~\cite{owaspmasvs}.

\begin{tableaucours}{L{4.2cm}L{8.0cm}}{%
    \tableauHeader{Domaine MASVS} &
    \tableauHeader{Ce qu'on vérifie}
}
Stockage des données & Secrets, jetons, données personnelles stockés en
clair sur l'appareil \\
Communication réseau & Chiffrement du canal, contrôle du certificat \\
Authentification & Jetons, biométrie, verrouillage \\
Code et résistance & Secrets codés en dur, journaux bavards \\
\end{tableaucours}
\vspace{-0.8em}
\begin{center}{\small\itshape Tableau~16.2. Domaines du MASVS}\end{center}

% ============================================================================
\section{L'Opération MEFOU}
\label{sec:c16-mefou}
% ============================================================================

\begin{boxscenario}{Opération MEFOU. Six constats sur le mobile}
\textbf{Constatations} (juillet 2026). (1)~L'application stocke le jeton de
session en clair sur l'appareil. (2)~Le service ne vérifie pas la propriété
de l'objet demandé~: un colis d'autrui est accessible par son numéro.
(3)~Une fonction d'administration est appelable par un compte ordinaire.
(4)~Le canal accepte un certificat non vérifié. (5)~Une clé de service est
codée en dur dans l'application. (6)~Les journaux de l'application contiennent
des données personnelles.

\medskip
\textbf{Analyse} (confiance élevée). Les constats (2) et (3), des défauts
d'accès, sont les plus graves~; ils préfigurent le constat T-05 du test
principal. L'application a été corrigée avant octobre, sauf le service, dont
le défaut d'accès subsiste et devient T-05.
\end{boxscenario}

\begin{boxmetier}{Pentesteur et RSSI}
\textbf{Tester le service, pas l'écran.} Pour une application mobile, le
pentesteur intercepte le trafic et rejoue les appels sans l'application~; le
RSSI retient que corriger l'application ne corrige pas le service, et que le
service doit refaire tous les contrôles.
\end{boxmetier}

\separateur

\begin{boxresume}
\textbf{Ce qu'il faut retenir de ce chapitre~:}
\begin{itemize}
    \item \textbf{L'interface de programmation} porte la logique~; on
          l'interroge directement.
    \item \textbf{Le défaut d'accès au niveau de l'objet} est le premier
          risque des API.
    \item \textbf{L'application mobile est une façade}~; ses contrôles ne
          protègent pas le service.
    \item \textbf{Le MASVS} structure le test mobile.
    \item \textbf{MEFOU} a montré, dès juillet, le défaut d'accès devenu
          T-05.
\end{itemize}
\end{boxresume}

\begin{boxinfo}
\textbf{Points de contrôle de l'élève.} Je sais~: décrire une API pour le
test~; citer les risques du Top~10 des API~; expliquer pourquoi le mobile ne
protège pas le service~; nommer les domaines du MASVS~; relier MEFOU à T-05.
\end{boxinfo}

\bigskip

\begin{boxexo}[title={\ding{45}\quad Exercice~16.1. API \niveaubase}]
Pourquoi teste-t-on une interface de programmation directement, sans
l'application qui l'utilise~?
\end{boxexo}

\begin{boxexo}[title={\ding{45}\quad Exercice~16.2. Premier risque \niveaubase}]
Quel est le premier risque du Top~10 des API~? À quoi correspond-il dans le
chapitre précédent~?
\end{boxexo}

\begin{boxexo}[title={\ding{45}\quad Exercice~16.3. Clé cachée \niveaubase}]
Pourquoi une clé « cachée » dans une application mobile n'est-elle pas
secrète~?
\end{boxexo}

\begin{boxexo}[title={\ding{45}\quad Exercice~16.4. Façade \niveaumoyen}]
Un développeur affirme que l'application empêche l'accès au colis d'autrui.
Que lui répondez-vous~?
\end{boxexo}

\begin{boxexo}[title={\ding{45}\quad Exercice~16.5. MASVS \niveaumoyen}]
Pour chacun des six constats de MEFOU, donnez le domaine MASVS concerné.
\end{boxexo}

\begin{boxexo}[title={\ding{45}\quad Exercice~16.6. MEFOU vers T-05 \niveaumoyen}]
Expliquez pourquoi le défaut d'accès de MEFOU se retrouve en octobre comme
T-05.
\end{boxexo}

\begin{boxexo}[title={\ding{45}\quad Exercice~16.7. Interception \niveaumoyen}]
Décrivez, sans outil précis, comment on observe les appels d'une application
mobile lors d'un test autorisé.
\end{boxexo}

\begin{boxexo}[title={\ding{45}\quad Exercice~16.8. Correction \niveaumoyen}]
Rédigez la correction du constat « fonction d'administration appelable par
un compte ordinaire ».
\end{boxexo}

\begin{boxexo}[title={\ding{45}\quad Exercice~16.9. Plan API \niveauexpert}]
Établissez le plan de test de l'interface de programmation de LiZbiZ selon
le Top~10 des API.
\end{boxexo}

\begin{boxexo}[title={\ding{45}\quad Exercice~16.10. Note au comité \niveauexpert}]
Rédigez une note qui explique au comité pourquoi « l'application a été
corrigée » ne suffit pas.
\end{boxexo}

\bigskip

\begin{boxinfo}
\textbf{Pour aller plus loin.}
\begin{itemize}
    \item OWASP API Security Top~10~\cite{owaspapitop10}~; OWASP
          MASVS~\cite{owaspmasvs}.
\end{itemize}
\end{boxinfo}

\begin{boxinfo}
\textbf{Transition.} Les services de LiZbiZ migrent vers le nuage~: le
Chapitre~\ref{chap:nuage} traite du test des environnements en nuage et des
conteneurs.
\end{boxinfo}
